\documentclass[12pt,a4paper]{article}
\usepackage[top=2.5cm,bottom=2.5cm,left=2.5cm,right=2.5cm,headheight=15pt]{geometry}
\usepackage{amsmath,amssymb}
\usepackage{tikz}
\usetikzlibrary{shapes,arrows,arrows.meta,positioning,calc,backgrounds,fit,matrix}
\usepackage{booktabs,array,longtable,tabularx,multirow}
\usepackage{xcolor}
\usepackage{enumitem}
\usepackage{fancyhdr}
\usepackage{titlesec}
\usepackage{mdframed}
\usepackage{tcolorbox}
\tcbuselibrary{skins,breakable}
\usepackage{hyperref}

%% ---- Colors ----
\definecolor{folblue}{RGB}{0,82,165}
\definecolor{lightblue}{RGB}{220,235,255}
\definecolor{darkorange}{RGB}{180,70,0}
\definecolor{lightorange}{RGB}{255,235,210}
\definecolor{darkgreen}{RGB}{20,110,20}
\definecolor{lightgreen}{RGB}{200,240,205}
\definecolor{darkred}{RGB}{150,20,20}
\definecolor{lightred}{RGB}{255,210,210}
\definecolor{lightgray}{RGB}{242,242,242}

%% ---- tcolorbox styles ----
\newtcolorbox{qbox}[2][]{enhanced,breakable,
  colback=lightorange,colframe=darkorange,arc=4pt,
  title={\textbf{#2}},fonttitle=\bfseries\color{white},
  attach boxed title to top left={yshift=-2mm,xshift=4mm},
  boxed title style={colback=darkorange,arc=3pt},#1}

\newtcolorbox{solbox}[2][]{enhanced,breakable,
  colback=lightblue,colframe=folblue,arc=4pt,
  title={\textbf{#2}},fonttitle=\bfseries\color{white},
  attach boxed title to top left={yshift=-2mm,xshift=4mm},
  boxed title style={colback=folblue,arc=3pt},#1}

\newtcolorbox{keybox}[2][]{enhanced,breakable,
  colback=lightgreen,colframe=darkgreen,arc=4pt,
  title={\textbf{#2}},fonttitle=\bfseries\color{white},
  attach boxed title to top left={yshift=-2mm,xshift=4mm},
  boxed title style={colback=darkgreen,arc=3pt},#1}

\newtcolorbox{algobox}[2][]{enhanced,breakable,
  colback=lightgray,colframe=gray!60,arc=4pt,
  title={\textbf{#2}},fonttitle=\bfseries,
  attach boxed title to top left={yshift=-2mm,xshift=4mm},
  boxed title style={colback=gray!60,arc=3pt},#1}

\newtcolorbox{warnbox}[2][]{enhanced,breakable,
  colback=lightred,colframe=darkred,arc=4pt,
  title={\textbf{#2}},fonttitle=\bfseries\color{white},
  attach boxed title to top left={yshift=-2mm,xshift=4mm},
  boxed title style={colback=darkred,arc=3pt},#1}

%% ---- Page style ----
\pagestyle{fancy}\fancyhf{}
\lhead{\color{folblue}\textbf{Classical Ciphers --- CSE 717 InfoSec}}
\rhead{\color{folblue}Exam: 17 Jun 2026}
\cfoot{--- \thepage\ ---}
\renewcommand{\headrulewidth}{0.4pt}

%% ---- Section style ----
\titleformat{\section}{\large\bfseries\color{folblue}}{}{0em}{}[\titlerule]
\titleformat{\subsection}{\normalsize\bfseries\color{darkorange}}{}{0em}{}

\begin{document}

\begin{center}
{\LARGE\bfseries\color{folblue} CSE 717 --- Information Security}\\[4pt]
{\large\bfseries Topic 1: Classical Ciphers}\\[2pt]
{\large Complete Past Paper Solutions: 2020--2024}\\[6pt]
{\normalsize Source: \textit{LcMaterial\#1\_InfoSecuritySP26.pdf} (59 pages, read in full)}\\[2pt]
{\normalsize\color{darkred}\textbf{Every classical-cipher question 2020--2024 included. Zero omissions.}}
\end{center}

\tableofcontents
\newpage

%% ==========================================
\section*{Quick Reference --- Algorithms}
%% ==========================================

\begin{keybox}{The Recurring Playfair Matrix --- MEMORIZE THIS}
2020, 2021, AND 2023 all GIVE the exact same $5\times5$ matrix (I/J share a cell):
\[
\begin{array}{|c|c|c|c|c|}
\hline
M & F & H & I/J & K \\\hline
U & N & O & P & Q \\\hline
Z & V & W & X & Y \\\hline
E & L & A & R & G \\\hline
D & S & T & B & C \\\hline
\end{array}
\]
High probability it recurs in 2026. If a matrix is given in the exam, you only need the encryption \emph{procedure} below --- do not waste time building a matrix from a keyword unless explicitly asked.
\end{keybox}

\begin{algobox}{Caesar Cipher}
\textbf{Encryption:} $C = (P + k) \bmod 26$ \quad \textbf{Decryption:} $P = (C - k) \bmod 26$,
where each letter is mapped A=0, B=1, \dots, Z=25.
General affine form sometimes used: $C=(aP+b)\bmod 26$ (with $a$ coprime to 26).
\end{algobox}

\begin{algobox}{Playfair Cipher --- Encryption Procedure}
\begin{enumerate}[noitemsep]
  \item Remove spaces/punctuation, uppercase the message; treat I and J as the same letter.
  \item Split into digraphs (pairs) left to right:
    \begin{itemize}[noitemsep]
      \item If both letters of a pair are identical, insert \texttt{X} between them and re-pair (the second of the pair shifts forward).
      \item If the message has odd length, append \texttt{X} to the final single letter.
    \end{itemize}
  \item For each digraph $(L_1, L_2)$ locate both letters in the matrix:
    \begin{itemize}[noitemsep]
      \item \textbf{Same row} $\to$ replace each letter with the letter immediately to its \emph{right} (wrap around row).
      \item \textbf{Same column} $\to$ replace each letter with the letter immediately \emph{below} it (wrap around column).
      \item \textbf{Rectangle} (different row \& column) $\to$ replace each letter with the letter in its \emph{own row} but the \emph{other letter's column}.
    \end{itemize}
\end{enumerate}
\end{algobox}

\begin{algobox}{Vigen\`ere Cipher}
Repeat the keyword to match the plaintext length. For position $i$:
\[ C_i = (P_i + K_i) \bmod 26 \] \quad (decryption: $P_i = (C_i - K_i)\bmod 26$)
\end{algobox}

\begin{algobox}{One-Time Pad (OTP)}
$C_i = (P_i + K_i) \bmod 26$ where $K$ is a truly random key, \textbf{as long as the message}, used \textbf{exactly once}, and kept secret. Provides Shannon \emph{perfect secrecy} --- ciphertext reveals zero information about plaintext --- but ONLY if all four conditions hold.
\end{algobox}

\begin{algobox}{Hill Cipher (2$\times$2, not yet in 2020--2024 papers but in Lc\#1)}
$C = (K \times P) \bmod 26$, where $P$ is a $2\times1$ plaintext vector and $K$ is an invertible $2\times2$ key matrix mod 26.
Decryption: $P = (K^{-1} \times C)\bmod 26$, where $K^{-1}=(\det K)^{-1}\begin{pmatrix}d&-b\\-c&a\end{pmatrix}\bmod 26$ for $K=\begin{pmatrix}a&b\\c&d\end{pmatrix}$.
\end{algobox}

\newpage
%% ==========================================
\section{2024 --- Q2(a,b): Caesar Encrypt + Decrypt (3 marks)}
%% ==========================================

\begin{qbox}{Question --- 2024 Q2}
\begin{enumerate}[label=(\alph*)]
  \item Encrypt the message ``CRYPTO'' first with a Caesar cipher shift of 5. \hfill [1.5]
  \item A message was encrypted using a Caesar cipher with the encryption function $f(p)=(p+4)\bmod 26$. The ciphertext received is ``ONXNG''. Determine the original plaintext. \hfill [1.5]
\end{enumerate}
\end{qbox}

\subsection{Part (a): Encrypt ``CRYPTO'', shift $k=5$}
\begin{solbox}{Caesar Encryption --- 2024 Q2(a)}
$C=(P+5)\bmod 26$, letter-by-letter:
\begin{center}
\renewcommand{\arraystretch}{1.3}
\begin{tabular}{c|c|c|c|c}
\toprule
Plain & Numeric $P$ & $P+5$ & $\bmod 26$ & Cipher \\
\midrule
C & 2  & 7  & 7  & H \\
R & 17 & 22 & 22 & W \\
Y & 24 & 29 & 3  & D \\
P & 15 & 20 & 20 & U \\
T & 19 & 24 & 24 & Y \\
O & 14 & 19 & 19 & T \\
\bottomrule
\end{tabular}
\end{center}
\textbf{Ciphertext: \texttt{HWDUYT}}
\end{solbox}

\subsection{Part (b): Decrypt ``ONXNG'', $f(p)=(p+4)\bmod 26$}
\begin{solbox}{Caesar Decryption --- 2024 Q2(b)}
Decryption is the inverse: $P=(C-4)\bmod 26$.
\begin{center}
\renewcommand{\arraystretch}{1.3}
\begin{tabular}{c|c|c|c|c}
\toprule
Cipher & Numeric $C$ & $C-4$ & $\bmod 26$ & Plain \\
\midrule
O & 14 & 10 & 10 & K \\
N & 13 & 9  & 9  & J \\
X & 23 & 19 & 19 & T \\
N & 13 & 9  & 9  & J \\
G & 6  & 2  & 2  & C \\
\bottomrule
\end{tabular}
\end{center}
\textbf{Plaintext: \texttt{KJTJC}}

\vspace{0.3em}
\noindent\textit{Note: the question only asks you to apply the inverse function --- the result need not be a dictionary word. Show the per-letter table; that is where the marks are.}
\end{solbox}

\newpage
%% ==========================================
\section{2023 --- Q1(a): Playfair Encryption with Given Matrix (4 marks)}
%% ==========================================

\begin{qbox}{Question --- 2023 Q1(a)}
Using the following Playfair matrix, encrypt the message: \textbf{``See you at CSECU''} \hfill [4]
\[
\begin{array}{|c|c|c|c|c|}
\hline
M & F & H & I/J & K \\\hline
U & N & O & P & Q \\\hline
Z & V & W & X & Y \\\hline
E & L & A & R & G \\\hline
D & S & T & B & C \\\hline
\end{array}
\]
\end{qbox}

\begin{solbox}{Playfair Encryption --- 2023 Q1(a)}
\textbf{Step 1 --- prepare the message.} Remove spaces, uppercase:
\[ \texttt{SEEYOUATCSECU} \quad (13 \text{ letters}) \]

\textbf{Step 2 --- form digraphs.} Pair sequentially; for a repeated-letter pair insert \texttt{X}; pad an odd final letter with \texttt{X}:
\[ \texttt{SE\ EY\ OU\ AT\ CS\ EC\ UX} \]

\textbf{Step 3 --- apply the row/column/rectangle rule to each digraph:}

\begin{center}
\renewcommand{\arraystretch}{1.3}
\begin{tabular}{c|l|c}
\toprule
Digraph & Rule applied & Cipher \\
\midrule
SE & Rectangle: S(row4,col1)$\to$col0=D; E(row3,col0)$\to$col1=L & DL \\
EY & Rectangle: E(row3,col0)$\to$col4=G; Y(row2,col4)$\to$col0=Z & GZ \\
OU & Same row (U N O P Q): O$\to$right=P; U$\to$right=N & PN \\
AT & Same column (H O W A T): A$\to$below=T; T$\to$below(wrap)=H & TH \\
CS & Same row (D S T B C): C$\to$right(wrap)=D; S$\to$right=T & DT \\
EC & Rectangle: E(row3,col0)$\to$col4=G; C(row4,col4)$\to$col0=D & GD \\
UX & Rectangle: U(row1,col0)$\to$col3=P; X(row2,col3)$\to$col0=Z & PZ \\
\bottomrule
\end{tabular}
\end{center}

\textbf{Ciphertext: \texttt{DL GZ PN TH DT GD PZ} $\to$ \texttt{DLGZPNTHDTGDPZ}}
\end{solbox}

\begin{warnbox}{Common mistakes}
\begin{itemize}[noitemsep]
  \item Forgetting to recombine ``See you at'' $\to$ \texttt{SEEYOUAT} \emph{without spaces} before pairing --- spaces are NOT encrypted as letters.
  \item Mixing up the rectangle rule direction: each letter takes \emph{its own row}, \emph{the other letter's column} --- not the reverse.
  \item Forgetting the wrap-around on same-row/same-column rules (e.g.\ rightmost column wraps to leftmost, bottom row wraps to top).
\end{itemize}
\end{warnbox}

\newpage
%% ==========================================
\section{2022 --- Q2(b): Caesar / OTP Cryptanalysis (conceptual, part of 10-mark Q2)}
%% ==========================================

\begin{qbox}{Question --- 2022 Q2}
\begin{enumerate}[label=(\alph*)]
  \item What are the essential ingredients of a symmetric cipher? \hfill [3]
  \item How can cryptanalysis possibly work for Caesar cipher and one-time pad? \hfill [3]
  \item What are the principal elements of a public-key cryptosystem? What are the roles of the public and the private key? \hfill [4]
\end{enumerate}
\end{qbox}

\subsection{Part (b): Cryptanalysis of Caesar Cipher and OTP}
\begin{solbox}{Conceptual Answer --- 2022 Q2(b)}
\textbf{Caesar cipher --- trivially broken:}
\begin{itemize}[noitemsep]
  \item \textbf{Brute force / exhaustive key search:} only 25 possible keys ($k=1..25$) --- try all shifts, the one producing readable text is correct. Computationally negligible.
  \item \textbf{Frequency analysis:} the shift preserves the relative frequency distribution of letters (e.g.\ \texttt{E} remains the most frequent letter, just relabelled). Comparing ciphertext letter frequencies to known English letter frequencies reveals the shift $k$ directly.
\end{itemize}

\textbf{One-time pad --- cryptanalysis only works if misused:}
\begin{itemize}[noitemsep]
  \item If used \emph{correctly} (truly random key, as long as the message, used once, kept secret) OTP gives \textbf{perfect secrecy} --- the ciphertext is statistically independent of the plaintext, so \emph{no} amount of computation or ciphertext can reveal any information. Cryptanalysis is mathematically impossible.
  \item Cryptanalysis ``works'' only when one of the conditions is violated:
  \begin{itemize}[noitemsep]
    \item \textbf{Key reuse} $\to$ XOR-ing two ciphertexts cancels the key, leaking $P_1 \oplus P_2$ (crib-dragging attacks).
    \item \textbf{Non-random key} (e.g.\ generated by a weak PRNG or derived from text) $\to$ key itself becomes analyzable.
    \item \textbf{Key shorter than message} (i.e.\ effectively a Vigen\`ere) $\to$ periodicity allows frequency analysis per key-position.
  \end{itemize}
\end{itemize}

\textbf{One-line answer:} Caesar falls to brute force/frequency analysis because the keyspace and structure are tiny; OTP cryptanalysis only succeeds when the ``one-time'', ``random'', or ``secret'' guarantees are broken.
\end{solbox}

\newpage
%% ==========================================
\section{2021 --- A2: OTP Discussion + Playfair (4.75 marks); A4(iv): Caesar Program (2 marks)}
%% ==========================================

\begin{qbox}{Question --- 2021 A2 \& A4(iv)}
\textbf{A-2.}
\begin{enumerate}[label=(\alph*)]
  \item[(a)] (i) The one-time pad offers complete security --- do you agree on this? Defend your answer. (ii) What are the fundamental difficulties of the one-time pad? \hfill [2.75]
  \item[(b)] Using the given Playfair matrix (same as above), encrypt the message: ``See you at CSE department''. \hfill [$\approx$2]
\end{enumerate}
\textbf{A-4(iv).} Write a program that can encrypt and decrypt using the general Caesar cipher. \hfill [2]
\end{qbox}

\subsection{Part A2(a): OTP --- ``Complete Security'' Discussion}
\begin{solbox}{OTP Conceptual Defence --- 2021 A2(a)}
\textbf{(i) Do you agree?} \textbf{Agree, but only under strict conditions.} Shannon proved OTP achieves \emph{perfect secrecy}: $H(P\mid C)=H(P)$ --- the ciphertext gives an attacker zero information about the plaintext, \emph{provided}:
\begin{enumerate}[noitemsep]
  \item the key is \textbf{truly random} (not pseudo-random),
  \item the key is \textbf{at least as long as} the message,
  \item the key is used \textbf{exactly once} and never reused, and
  \item the key is kept \textbf{completely secret} between sender and receiver.
\end{enumerate}
If even one condition fails, the security guarantee collapses entirely.

\textbf{(ii) Fundamental difficulties (why OTP is impractical):}
\begin{itemize}[noitemsep]
  \item \textbf{Key generation:} producing truly random bits at the required volume is hard.
  \item \textbf{Key distribution:} a key as long as the message must be shared \emph{in advance, securely} --- as hard as securing the message itself.
  \item \textbf{Key storage:} both parties must store, protect, and synchronize huge key material.
  \item \textbf{No reuse:} every new message needs fresh key material --- keys cannot be cached or compressed.
  \item \textbf{Scalability:} impractical for high-volume or long-running communication (e.g.\ internet traffic).
\end{itemize}
\end{solbox}

\subsection{Part A2(b): Playfair --- ``See you at CSE department''}
\begin{solbox}{Playfair Encryption --- 2021 A2(b)}
\textbf{Step 1.} Remove spaces, uppercase: \texttt{SEEYOUATCSEDEPARTMENT} (21 letters).

\textbf{Step 2.} Digraphs (no repeated-letter pairs occur; final letter \texttt{T} is alone $\to$ pad with \texttt{X}):
\[ \texttt{SE\ EY\ OU\ AT\ CS\ ED\ EP\ AR\ TM\ EN\ TX} \]

\textbf{Step 3.} Apply rules (first five digraphs reuse the 2023 Q1(a) results):
\begin{center}
\renewcommand{\arraystretch}{1.3}
\begin{tabular}{c|l|c}
\toprule
Digraph & Rule applied & Cipher \\
\midrule
SE & (as 2023 Q1a) Rectangle & DL \\
EY & (as 2023 Q1a) Rectangle & GZ \\
OU & (as 2023 Q1a) Same row & PN \\
AT & (as 2023 Q1a) Same column & TH \\
CS & (as 2023 Q1a) Same row & DT \\
ED & Same column (M U Z E D): E$\to$below=D; D$\to$below(wrap)=M & DM \\
EP & Rectangle: E(row3,col0)$\to$col3=R; P(row1,col3)$\to$col0=U & RU \\
AR & Same row (E L A R G): A$\to$right=R; R$\to$right=G & RG \\
TM & Rectangle: T(row4,col2)$\to$col0=D; M(row0,col0)$\to$col2=H & DH \\
EN & Rectangle: E(row3,col0)$\to$col1=L; N(row1,col1)$\to$col0=U & LU \\
TX & Rectangle: T(row4,col2)$\to$col3=B; X(row2,col3)$\to$col2=W & BW \\
\bottomrule
\end{tabular}
\end{center}

\textbf{Ciphertext:} \texttt{DL GZ PN TH DT DM RU RG DH LU BW} $\to$ \texttt{DLGZPNTHDTDMRURGDHLUBW}
\end{solbox}

\subsection{Part A4(iv): Caesar Cipher Program}
\begin{solbox}{Python Program --- General Caesar Encrypt/Decrypt --- 2021 A4(iv)}
\begin{verbatim}
def caesar_encrypt(text, k):
    result = ""
    for ch in text.upper():
        if ch.isalpha():
            result += chr((ord(ch) - ord('A') + k) % 26 + ord('A'))
        else:
            result += ch
    return result

def caesar_decrypt(text, k):
    return caesar_encrypt(text, -k)

# Example
print(caesar_encrypt("CRYPTO", 5))   # -> HWDUYT
print(caesar_decrypt("HWDUYT", 5))   # -> CRYPTO
\end{verbatim}
\textit{Exam tip:} if asked to ``write a program'', pseudocode or Python both earn marks --- the examiner is checking you understand the modular-arithmetic loop, not syntax perfection.
\end{solbox}

\newpage
%% ==========================================
\section{2020 --- Q1(a,b): Playfair Encryption (4 marks) + Vigen\`ere Encryption (4.75 marks)}
%% ==========================================

\begin{qbox}{Question --- 2020 Q1}
\begin{enumerate}[label=(\alph*)]
  \item Using the given Playfair matrix (same matrix as above), encrypt this message: ``Must see you over CU playground. Coming now.'' \hfill [4]
  \item Using the Vigen\`ere cipher, encrypt the word ``explanation'' using the key ``cse''. \hfill [4.75]
\end{enumerate}
\end{qbox}

\subsection{Part (a): Playfair --- ``Must see you over CU playground. Coming now.''}
\begin{solbox}{Playfair Encryption --- 2020 Q1(a)}
\textbf{Step 1.} Strip punctuation/spaces, uppercase:
\[ \texttt{MUSTSEEYOUOVERCUPLAYGROUNDCOMINGNOW} \quad (35 \text{ letters}) \]

\textbf{Step 2.} No adjacent-pair repeats occur; pair sequentially and pad the trailing single letter \texttt{W} with \texttt{X}:
\[ \texttt{MU\,ST\,SE\,EY\,OU\,OV\,ER\,CU\,PL\,AY\,GR\,OU\,ND\,CO\,MI\,NG\,NO\,WX} \quad (18\text{ digraphs}) \]

\textbf{Step 3.} Apply the row/column/rectangle rule:
\begin{center}
\renewcommand{\arraystretch}{1.25}
\begin{tabular}{c|l|c}
\toprule
Digraph & Rule applied & Cipher \\
\midrule
MU & Same col (M U Z E D): M$\to$below=U; U$\to$below=Z & UZ \\
ST & Same row (D S T B C): S$\to$right=T; T$\to$right=B & TB \\
SE & Rectangle (as before) & DL \\
EY & Rectangle (as before) & GZ \\
OU & Same row (as before) & PN \\
OV & Rectangle: O(row1,col2)$\to$col1=N; V(row2,col1)$\to$col2=W & NW \\
ER & Same row (E L A R G): E$\to$right=L; R$\to$right=G & LG \\
CU & Rectangle: C(row4,col4)$\to$col0=D; U(row1,col0)$\to$col4=Q & DQ \\
PL & Rectangle: P(row1,col3)$\to$col1=N; L(row3,col1)$\to$col3=R & NR \\
AY & Rectangle: A(row3,col2)$\to$col4=G; Y(row2,col4)$\to$col2=W & GW \\
GR & Same row (E L A R G): G$\to$right(wrap)=E; R$\to$right=G & EG \\
OU & Same row (as before) & PN \\
ND & Rectangle: N(row1,col1)$\to$col0=U; D(row4,col0)$\to$col1=S & US \\
CO & Rectangle: C(row4,col4)$\to$col2=T; O(row1,col2)$\to$col4=Q & TQ \\
MI & Same row (M F H I/J K): M$\to$right=F; I/J$\to$right=K & FK \\
NG & Rectangle: N(row1,col1)$\to$col4=Q; G(row3,col4)$\to$col1=L & QL \\
NO & Same row (U N O P Q): N$\to$right=O; O$\to$right=P & OP \\
WX & Same row (Z V W X Y): W$\to$right=X; X$\to$right=Y & XY \\
\bottomrule
\end{tabular}
\end{center}

\textbf{Ciphertext:} \texttt{UZTBDLGZPNNWLGDQNRGWEGPNUSTQFKQLOPXY}

\textit{Exam tip: with 35 letters this is the longest Playfair question seen 2020--2024 (4 marks). Work systematically left-to-right in a table like above --- partial credit is given per-digraph, so do not skip steps even under time pressure.}
\end{solbox}

\subsection{Part (b): Vigen\`ere --- ``explanation'', key ``cse''}
\begin{solbox}{Vigen\`ere Encryption --- 2020 Q1(b)}
Plaintext: \texttt{EXPLANATION} (11 letters). Repeat key \texttt{CSE} to length 11: \texttt{CSECSECSECS}.
$C_i=(P_i+K_i)\bmod 26$:

\begin{center}
\renewcommand{\arraystretch}{1.3}
\begin{tabular}{c|c|c|c|c|c|c}
\toprule
Plain & $P_i$ & Key & $K_i$ & $P_i+K_i$ & $\bmod 26$ & Cipher \\
\midrule
E & 4  & C & 2  & 6  & 6  & G \\
X & 23 & S & 18 & 41 & 15 & P \\
P & 15 & E & 4  & 19 & 19 & T \\
L & 11 & C & 2  & 13 & 13 & N \\
A & 0  & S & 18 & 18 & 18 & S \\
N & 13 & E & 4  & 17 & 17 & R \\
A & 0  & C & 2  & 2  & 2  & C \\
T & 19 & S & 18 & 37 & 11 & L \\
I & 8  & E & 4  & 12 & 12 & M \\
O & 14 & C & 2  & 16 & 16 & Q \\
N & 13 & S & 18 & 31 & 5  & F \\
\bottomrule
\end{tabular}
\end{center}

\textbf{Ciphertext: \texttt{GPTNSRCLMQF}}
\end{solbox}

\newpage
%% ==========================================
\section*{Exam Strategy Summary}
%% ==========================================

\begin{keybox}{High-Yield Checklist for Classical Ciphers (5/5 years, 2--9 marks)}
\begin{enumerate}[noitemsep]
  \item \textbf{Caesar encrypt + decrypt} with arbitrary $f(p)=(ap+b)\bmod 26$ --- both directions, every year since 2024 and 2021's program question. Practice the per-letter table format.
  \item \textbf{Playfair with a GIVEN matrix} --- memorize the recurring matrix (2020/2021/2023). Drill: prepare message (strip spaces/punctuation, handle repeats with X, pad odd length), then row/column/rectangle rules. Practice on sentences \emph{with spaces} since that's the exam style.
  \item \textbf{Vigen\`ere} (2020) --- short keyword repeated, $C_i=(P_i+K_i)\bmod 26$. Mechanical, fast marks.
  \item \textbf{OTP conceptual} (2021, 2022) --- be ready to argue ``perfect secrecy under 4 strict conditions'' and list the practical difficulties / cryptanalysis-only-on-misuse points.
  \item \textbf{Hill Cipher} (Lc\#1 only, no past-paper appearance) --- light pass: know $C=(KP)\bmod 26$ and be able to reproduce the worked ``HI''$\to$``TC'' example if time allows, but do not over-invest.
\end{enumerate}
\end{keybox}

\end{document}
